\documentclass[10pt,a4paper]{amsart}
\usepackage[margin=19mm]{geometry}
\usepackage{amsmath,amssymb,mathtools}
\usepackage[scr=rsfs,cal=euler]{mathalfa}
\usepackage{xfrac}
\usepackage[unicode,hidelinks,bookmarks=false]{hyperref}
\newcommand{\ie}{i.e.\ }
\newcommand{\cf}{cf.\ }
\newcommand{\resp}{respectively\ }
\newcommand{\Subsection}{\subsection}
\providecommand{\nobreakdash}{\nobreak\hbox{-}}
\setlength{\emergencystretch}{2em}
\multlinegap0pt
\usepackage{csquotes}
\usepackage{amssymb}
\usepackage{xcolor}
\newtheorem{lemma}{Lemma}
\newcommand{\Ad}{\operatorname{Ad}}
\newcommand{\psum}{\sideset{}{^*}\sum}
\title{On Ramanujan Primes for Hecke-Maass cusp forms}
\author{Tinghao Huang, Shifan Zhao}
\date{Source: arXiv:2605.09807v1 (CC BY 4.0)}
\begin{document}
\maketitle

\noindent\emph{Source and licence.} On Ramanujan Primes for Hecke-Maass cusp forms. Version arXiv:2605.09807v1 (CC BY 4.0), distributed by arXiv under the Creative Commons Attribution 4.0 International licence, http://creativecommons.org/licenses/by/4.0/, as recorded in the arXiv licence metadata for this exact version and shown on the abstract page https://arxiv.org/abs/2605.09807v1. Full licence notice: \texttt{licenses/arxiv-2605.09807-CC-BY-4.0.txt}.\par\smallskip

\noindent\textit{Editorial scope.} Counted unit: the article's lemma UpperBoundforS(x) (source0.tex lines 286--292). The article's complete original proof of it is delivered with it (source0.tex lines 294--331, one \texttt{proof} environment of 38 lines). No mathematics is written, repaired or re-derived: the statement and the proof are the article's own lines, in the article's own notation, and no retained line is substituted or re-flowed.  Retained uncounted as the source-local context the proof needs: lines 110--112; lines 223--225; lines 232--234; lines 241--243.  Nothing was omitted from any retained span, so the package carries no omission record.  No retained line contains a citation, so the wrapper carries no bibliography and no bibliography entry.  No retained line contains a figure environment, an image-inclusion command or drawing code, so no image is delivered and the package declares no asset dependency.  Editorial formatting only, in addition: an AMS \texttt{amsart} preamble that restates the article's own declarations and macros used by the retained lines, namely the environments \texttt{lemma} on separate counters, together with the standard packages those lines need.  The retained environments print in this document as Lemma 1 in order of appearance, whereas the article prints the same items with the numbering of its own full document.  The complete native source of the article as distributed in the arXiv e-print for this exact version is bundled unchanged as \texttt{sources/200262/source0.tex}.
\begin{equation}\label{KimSarnakboundforphi}
     |\lambda_{\phi}(p)| \leq p^{\frac{7}{64}} + p^{-\frac{7}{64}}.
\end{equation}

\begin{equation}\label{EulerProductofAd}
    L(s,\Ad(\phi)) = \prod_{p\nmid N}\bigg[1 - \frac{A_{\phi}(p)}{p^s} + \frac{A_{\phi}(p)}{p^{2s}} - \frac{1}{p^{3s}}\bigg]^{-1}
\end{equation}

\begin{equation}\label{convexitybound}
    L(\sigma + it, \Ad(\phi)) \ll_{\epsilon} \big[Q_{\phi}\cdot (1 + |t|)^3\big]^{\frac{1-\sigma}{2} + \epsilon}.
\end{equation}

\begin{equation}\label{>3atp^2}
    \lambda_{\phi_i}(p^2)\overline{\chi_{\phi_i}(p)} = |\lambda_{\phi_i}(p)|^2 - 1
\end{equation}

\begin{lemma}\label{UpperBoundforS(x)}
    We have for any $\frac{23}{32} + \epsilon< \sigma < 1$: 
    \[
        S(x) \ll_{\epsilon} x^{\sigma+\epsilon}\cdot Q_{\phi_1}^{\frac{1-\sigma}{2}+\epsilon}Q_{\phi_2}^{\frac{1-\sigma}{2}+\epsilon}
    \]
    where the implied constant depends only on $\epsilon > 0$.
\end{lemma}

\begin{proof}
    We define the following auxiliary Euler product:
    \begin{align*}
        G(s)  & := \prod_{p\nmid N_1N_2} \bigg[\bigg(1 - \frac{A_{\phi_1}(p)}{p^s} + \frac{A_{\phi_1}(p)}{p^{2s}} - \frac{1}{p^{3s}}\bigg) \bigg(1 - \frac{A_{\phi_2}(p)}{p^s} + \frac{A_{\phi_2}(p)}{p^{2s}} - \frac{1}{p^{3s}}\bigg)\bigg]\cdot \bigg[1 + \frac{B(p)}{p^s}\bigg] \\
        & \times \prod_{\substack{p\nmid N_1\\p|N_2}}\bigg(1 - \frac{A_{\phi_1}(p)}{p^s} + \frac{A_{\phi_1}(p)}{p^{2s}} - \frac{1}{p^{3s}}\bigg)\cdot 
        \prod_{\substack{p\nmid N_2\\p|N_1}}\bigg(1 - \frac{A_{\phi_2}(p)}{p^s} + \frac{A_{\phi_2}(p)}{p^{2s}} - \frac{1}{p^{3s}}\bigg)
    \end{align*}
    By \eqref{KimSarnakboundforphi} and \eqref{>3atp^2}, we have 
    \begin{equation}\label{KimSarnakBoudnforAp}
        |A_{\phi_i}(p)| \leq p^{\frac{7}{32}} + p^{\frac{-7}{32}} + 1.
    \end{equation}
    Then, upon expanding the factors and making use of \eqref{KimSarnakBoudnforAp}, we see that for $p\nmid N_1N_2$ we have
    \begin{align*}
        \bigg[\bigg(1 - \frac{A_{\phi_1}(p)}{p^s} + \frac{A_{\phi_1}(p)}{p^{2s}} - \frac{1}{p^{3s}}\bigg) \bigg(1 - \frac{A_{\phi_2}(p)}{p^s} + \frac{A_{\phi_2}(p)}{p^{2s}} - \frac{1}{p^{3s}}\bigg)\bigg]\cdot \bigg[1 + \frac{B(p)}{p^s}\bigg] 
        = 1 + O(p^{-2\sigma + \frac{7}{16}}),
    \end{align*}
    where $\sigma := \Re(s)$.
    We therefore conclude that for $\sigma > \frac{23}{32}+\epsilon$, $G(s)$ converges absolutely and we have
    \[
        |G(s)| \ll \tau(N_1)\tau(N_2)\cdot \prod_{p\nmid N_1N_2}(1 + O(p^{-1 - 2\epsilon})) \ll_{\epsilon} (N_1N_2)^{\epsilon},
    \]
    where the implied constant depends only on $\epsilon > 0$ and $\tau(\cdot)$ is the divisor function.
    
    Now, as for $\sigma > 1$ we have 
    \[
        L(s,\Ad(\phi_1))L(s,\Ad(\phi_2))G(s) = \psum_{\substack{n \geq 1 \\ (n,N_1N_2)=1}}\frac{B(n)}{n^s}
    \]
    by \eqref{EulerProductofAd}, for $c>1$ we could relate $S(x)$ to a line integral as
    \begin{align*}
        S(x) & = \frac{1}{2\pi i}\int_{(c)}L(s,\Ad(\phi_1))L(s,\Ad(\phi_2))G(s)\frac{x^s}{s^2}ds.
    \end{align*}
    As $L(s, \Ad(\phi_1)), L(s, \Ad(\phi_2))$ are continued and $G(s)$ is absolutely convergent for $\sigma > \frac{23}{32} + \epsilon$, we may shift the integral line left to any vertical line with real part $\sigma$ satisfying $\frac{23}{32} +\epsilon < \sigma <1$ and conclude
    \[
        S(x) \ll_{\epsilon} x^{\sigma+\epsilon}\cdot (Q_{\phi_1}Q_{\phi_2})^{\frac{1-\sigma}{2} + \epsilon}
    \]
    in the view of the convexity bounds for $\Ad(\phi_1)$ and $\Ad(\phi_2)$ as in \eqref{convexitybound}.
    
\end{proof}

\end{document}
